\مسئله{}
یک سیستم LTI زمان‌گسسته با پاسخ فرکانسی $H(e^{j\omega})$ و پاسخ ضربه متناظر $h[n]$ را در نظر بگیرید.

الف) ۳ نشانه زیر در مورد این سیستم داده شده است:
\begin{itemize}
    \item سیستم علی (Causal) است.
    \item $H(e^{j\omega}) = H^*(e^{-j\omega})$
    \item تبدیل فوریه زمان‌گسسته (DTFT) دنباله $h[n+1]$ حقیقی است.
\end{itemize}
با استفاده از این ۳ نشانه نشان دهید که پاسخ ضربه سیستم دارای مدت زمان محدود است.

ب) علاوه بر نشانه‌های قبلی، دو نشانه دیگر نیز داده شده است:
\begin{itemize}
    \item $\frac{1}{2\pi} \int_{\pi}^{\pi} H(e^{j\omega}) d\omega = 2$
    \item $H(e^{j\pi}) = 0$
\end{itemize}
آیا اطلاعات کافی برای تعیین منحصربه‌فرد سیستم وجود دارد؟ اگر پاسخ مثبت است، پاسخ ضربه $h[n]$ را تعیین کنید. در غیر این صورت، تا حد امکان ویژگی‌های دنباله $h[n]$ را مشخص کنید.

\textbf{بخش امتیازی:} مسئله را در حالتی حل کنید که نشانه چهارم به صورت انرژی بیان شده باشد:
\[
\frac{1}{2\pi} \int_{-\pi}^{\pi} |H(e^{j\omega})|^2 d\omega = 2
\]


\حل{}

\textbf{پاسخ بخش الف)} 
تفسیر نشانه‌های داده شده به شرح زیر است:
\begin{itemize}
    \item از نشانه (ii) داریم $H(e^{j\omega}) = H^*(e^{-j\omega})$ که نشان می‌دهد سیگنال حوزه زمان آن یعنی $h[n]$ کاملاً حقیقی است ($h[n] = h^*[n]$).
    \item از نشانه (i) به دلیل علیت سیستم، پاسخ ضربه برای زمان‌های منفی صفر است، یعنی:
    \[ h[n] = 0 \quad (\forall n < 0) \]
    \item از نشانه (iii) ، چون DTFT دنباله $h[n+1]$ حقیقی است و خود سیگنال نیز حقیقی می‌باشد، این دنباله در حوزه زمان باید زوج باشد:
    \[ h[n+1] = h[-n+1] \]
\end{itemize}
با توجه به شرط علیت، داریم $h[n] = 0$ برای $n < 0$، در نتیجه $h[n+1] = 0$ برای $n < -1$ خواهد بود. از طرفی چون $h[n+1]$ زوج است، صفر بودن آن در $n < -1$ ایجاب می‌کند که برای $n > 1$ نیز صفر باشد. 

بنابراین دنباله $h[n+1]$ تنها می‌تواند در مقادیر $n \in \{-1, 0, 1\}$ غیرصفر باشد که این یعنی پاسخ ضربه اصلی سیستم ($h[n]$) تنها در نقاط $n \in \{0, 1, 2\}$ مقدار دارد و طول آن محدود است.

\newpage

\textbf{پاسخ بخش ب)}
تفسیر دو نشانه جدید به صورت زیر است:
\begin{itemize}
    \item از نشانه (iv) مقدار پاسخ ضربه در مبدأ مشخص می‌شود:
    \[ h[0] = \frac{1}{2\pi} \int_{-\pi}^{\pi} H(e^{j\omega}) d\omega = 2 \]
    چون طبق بخش قبل $h[n+1]$ زوج بود، داریم $h[1+1] = h[-1+1]$ که نتیجه می‌دهد:
    \[ h[2] = h[0] = 2 \]
    \item از نشانه (v) رابطه فرکانسی در $\omega = \pi$ را باز می‌کنیم:
    \[
    \begin{aligned}
    H(e^{j\pi}) &= \sum_{n=-\infty}^{\infty} h[n] e^{-j\pi n} \\
    &= h[0]e^0 + h[1]e^{-j\pi} + h[2]e^{-j2\pi} \\
    &= h[0] - h[1] + h[2] = 0
    \end{aligned}
    \]
\end{itemize}
با جایگذاری مقادیر به دست آمده در رابطه نهایی:
\[ 2 - h[1] + 2 = 0 \implies h[1] = 4 \]
بنابراین اطلاعات کافی است و سیستم به طور منحصربه‌فرد با پاسخ ضربه زیر مشخص می‌شود:
\[ h[n] = 2\delta[n] + 4\delta[n-1] + 2\delta[n-2] \]

\textbf{پاسخ بخش امتیازی}
با جایگزین شدن شرط چهارم با رابطه انرژی، طبق قضیه پارسوال داریم:
\[ \frac{1}{2\pi} \int_{-\pi}^{\pi} |H(e^{j\omega})|^2 d\omega = \sum_{n=-\infty}^{\infty} |h[n]|^2 = 2 \]
با توجه به اینکه سیگنال فقط در نقاط $n \in \{0, 1, 2\}$ غیرصفر است، رابطه فوق به صورت زیر ساده می‌شود:
\[ h[0]^2 + h[1]^2 + h[2]^2 = 2 \]
از بخش‌های قبل می‌دانیم که $h[2] = h[0]$ و همچنین از شرط $H(e^{j\pi}) = 0$ رابطه $h[1] = 2h[0]$ برقرار است. با جایگذاری این مقادیر در معادله انرژی خواهیم داشت:
\[
\begin{aligned}
h[0]^2 + (2h[0])^2 + h[0]^2 &= 2 \\
&= h[0]^2 + 4h[0]^2 + h[0]^2 \\
&= 6h[0]^2 = 2
\end{aligned}
\]
در نتیجه مقدار $h[0]$ برابر است با:
\[ h[0]^2 = \frac{1}{3} \implies h[0] = \pm \frac{1}{\sqrt{3}} \]
بنابراین در این حالت دو سیستم منحصربه‌فرد با پاسخ‌های ضربه زیر قابل تعریف خواهند بود:
\[ 
h_1[n] = \frac{1}{\sqrt{3}}\delta[n] + \frac{2}{\sqrt{3}}\delta[n-1] + \frac{1}{\sqrt{3}}\delta[n-2]
\]
\[
h_2[n] = -\frac{1}{\sqrt{3}}\delta[n] - \frac{2}{\sqrt{3}}\delta[n-1] - \frac{1}{\sqrt{3}}\delta[n-2]
\]